\documentclass[10pt, letterpaper]{article}

% Packages:
\usepackage[
  ignoreheadfoot, % set margins without considering header and footer
  top=2 cm, % seperation between body and page edge from the top
  bottom=2 cm, % seperation between body and page edge from the bottom
  left=2 cm, % seperation between body and page edge from the left
  right=2 cm, % seperation between body and page edge from the right
  footskip=1.0 cm, % seperation between body and footer
  % showframe % for debugging
]{geometry} % for adjusting page geometry
\usepackage{titlesec} % for customizing section titles
\usepackage{tabularx} % for making tables with fixed width columns
\usepackage{array} % tabularx requires this
\usepackage[dvipsnames]{xcolor} % for coloring text
\definecolor{primaryColor}{RGB}{0, 0, 0} % define primary color
\usepackage{enumitem} % for customizing lists
\usepackage{fontawesome5} % for using icons
\usepackage{amsmath} % for math
\usepackage[
  pdftitle={Samuel Nandi's CV},
  pdfauthor={Samuel Nandi},
  colorlinks=true,
  urlcolor=primaryColor
]{hyperref} % for links, metadata and bookmarks
\usepackage[pscoord]{eso-pic} % for floating text on the page
\usepackage{calc} % for calculating lengths
\usepackage{bookmark} % for bookmarks
\usepackage{lastpage} % for getting the total number of pages
\usepackage{changepage} % for one column entries (adjustwidth environment)
  \usepackage{paracol} % for two and three column entries
  \usepackage{ifthen} % for conditional statements
  \usepackage{needspace} % for avoiding page brake right after the section title
  \usepackage{iftex} % check if engine is pdflatex, xetex or luatex

  % Ensure that generate pdf is machine readable/ATS parsable:
  \ifPDFTeX
  \input{glyphtounicode}
  \pdfgentounicode=1
  \usepackage[T1]{fontenc}
  \usepackage[utf8]{inputenc}
  \usepackage{lmodern}
  \fi

  \usepackage{charter}

  % Some settings:
  \raggedright
  \AtBeginEnvironment{adjustwidth}{\partopsep0pt} % remove space before adjustwidth environment
  \pagestyle{empty} % no header or footer
  \setcounter{secnumdepth}{0} % no section numbering
  \setlength{\parindent}{0pt} % no indentation
  \setlength{\topskip}{0pt} % no top skip
  \setlength{\columnsep}{0.15cm} % set column seperation
  \pagenumbering{gobble} % no page numbering

  \titleformat{\section}{\needspace{4\baselineskip}\bfseries\large}{}{0pt}{}[\vspace{1pt}\titlerule]

  \titlespacing{\section}{
    % left space:
      -1pt
  }{
    % top space:
      0.3 cm
  }{
    % bottom space:
      0.2 cm
  } % section title spacing

\renewcommand\labelitemi{$\vcenter{\hbox{\small$\bullet$}}$} % custom bullet points
\newenvironment{highlights}{
  \begin{itemize}[
    topsep=0.10 cm,
    parsep=0.10 cm,
    partopsep=0pt,
    itemsep=0pt,
    leftmargin=0 cm + 10pt
  ]
}{
  \end{itemize}
} % new environment for highlights


\newenvironment{highlightsforbulletentries}{
  \begin{itemize}[
    topsep=0.10 cm,
    parsep=0.10 cm,
    partopsep=0pt,
    itemsep=0pt,
    leftmargin=10pt
  ]
}{
  \end{itemize}
} % new environment for highlights for bullet entries

\newenvironment{onecolentry}{
  \begin{adjustwidth}{
    0 cm + 0.00001 cm
  }{
    0 cm + 0.00001 cm
  }
}{
  \end{adjustwidth}
} % new environment for one column entries

\newenvironment{twocolentry}[2][]{
  \onecolentry
    \def\secondColumn{#2}
  \setcolumnwidth{\fill, 4.5 cm}
  \begin{paracol}{2}
}{
  \switchcolumn \raggedleft \secondColumn
    \end{paracol}
  \endonecolentry
} % new environment for two column entries

\newenvironment{threecolentry}[3][]{
  \onecolentry
    \def\thirdColumn{#3}
  \setcolumnwidth{, \fill, 4.5 cm}
  \begin{paracol}{3}
  {\raggedright #2} \switchcolumn
}{
  \switchcolumn \raggedleft \thirdColumn
    \end{paracol}
  \endonecolentry
} % new environment for three column entries

\newenvironment{header}{
  \setlength{\topsep}{0pt}\par\kern\topsep\centering\linespread{1.5}
}{
  \par\kern\topsep
} % new environment for the header

\newcommand{\placelastupdatedtext}{% \placetextbox{<horizontal pos>}{<vertical pos>}{<stuff>}
  \AddToShipoutPictureFG*{% Add <stuff> to current page foreground
    \put(
        \LenToUnit{\paperwidth-2 cm-0 cm+0.05cm},
        \LenToUnit{\paperheight-1.0 cm}
        ){\vtop{{\null}\makebox[0pt][c]{
      \small\color{gray}\textit{Last updated in September 2024}\hspace{\widthof{Last updated in September 2024}}
    }}}%
  }%
}%

% save the original href command in a new command:
\let\hrefWithoutArrow\href

% new command for external links:

\begin{document}
\newcommand{\AND}{\unskip
  \cleaders\copy\ANDbox\hskip\wd\ANDbox
    \ignorespaces
}
\newsavebox\ANDbox
\sbox\ANDbox{$|$}

\begin{header}
\fontsize{25 pt}{25 pt}\selectfont Samuel Nandi

\vspace{5 pt}

\normalsize
\mbox{\hrefWithoutArrow{https://github.com/ssnnee}{Github}}%
\kern 5.0 pt%
\AND%
\kern 5.0 pt%
\mbox{\hrefWithoutArrow{mailto:nandisamuelsne@proton.me}{Mail}}%
\kern 5.0 pt%
\AND%
\kern 5.0 pt%
\mbox{\hrefWithoutArrow{https://linkedin.com/in/samuel-nandi}{Linkedin}}%
\kern 5.0 pt%
\AND%
\kern 5.0 pt%
\mbox{Brazzaville, Republic of Congo}%
\kern 5.0 pt%
\AND%
\kern 5.0 pt%
\mbox{\hrefWithoutArrow{tel:+242-04-439-2033}{+242 04 439 2833}}%
\end{header}

\section{About Me}
\begin{onecolentry}
  I am a software developer passionate about building solutions for everyday
  problems. Currently, I primarily work with high-level languages and frameworks,
  but I have interest in learning both modern and legacy technologies.
  In the future, I would like to contribute to or build low-level software and
  core tools that empower other developers to create their own solutions.
\end{onecolentry}

\section{Education}

\begin{twocolentry}{
    Feb 2023 – Oct 2024
}
\textbf{\href{https://www.alxafrica.com/}{ALX Software Engineering Program}} | \href{https://drive.google.com/file/d/1xa94zxY_QAznumHXPoMkZT9yH5yQwSzU/view?usp=sharing}{Certificate of Software Engineering}
\end{twocolentry}

% \begin{onecolentry}
% \begin{highlights}
% \item Intensive, project-based curriculum covering full-stack development and system engineering
% \end{highlights}
% \end{onecolentry}

\vspace{0.10 cm}
\begin{twocolentry}{
    Oct 2023 – Jul 2024
}
\textbf{\href{https://esgae.org/}{ESGAE}} | \href{https://drive.google.com/file/d/1XByQTQBimuyxj4Y3iiMretpPVR5-1aT8/view?usp=sharing}{Bachelor of Software Engineering}
\end{twocolentry}

\vspace{0.10 cm}
% \begin{onecolentry}
% \begin{highlights}
%   \item Focused on software architecture, design patterns, and enterprise application development
% \end{highlights}
% \end{onecolentry}

\begin{twocolentry}{
    Oct 2021 – Jul 2023
}
\textbf{\href{https://esgae.org/}{ESGAE}} | \href{https://drive.google.com/file/d/1yH2_KIsGoWv-wQOauQIn6xb9ZHOo6eEs/view?usp=sharing}{Associate of Analysis and Programming}
\end{twocolentry}

\vspace{0.10 cm}
% \begin{onecolentry}
% \begin{highlights}
% \item Coursework: C\#, Java, SQL, Web Development (HTML, CSS, JavaScript, PHP), UML, Merise
% \item Fundamentals of algorithms, database design, and software modeling
% \end{highlights}
% \end{onecolentry}

\section{Projects}

\begin{twocolentry}{
  \href{https://github.com/ssnnee/cellotree_web}{Github}
}
\textbf{CelloTree - Family Tree Web Application}
\end{twocolentry}

\vspace{0.10 cm}
\begin{onecolentry}
\begin{highlights}
\item A full-stack web application for creating and sharing family trees
\item Technologies: TypeScript, Next.js, Prisma, tRPC, PostgreSQL, Docker
\end{highlights}
\end{onecolentry}

\begin{twocolentry}{
  \href{https://github.com/ssnnee/kolandela}{Github}
}
\textbf{Kolandéla - Personal Finance Mobile App}
\end{twocolentry}

\vspace{0.10 cm}
\begin{onecolentry}
\begin{highlights}
\item A cross-platform mobile application for personal finance management
% \item Developed a RESTful API backend using Go with clean architecture principles
% \item Implemented efficient data synchronization between local SQLite and server
% \item Set up HAProxy for load balancing and improved application reliability
\item Technologies: TypeScript, Go, React Native, SQLite
\end{highlights}
\end{onecolentry}

\section{Technical Skills}

\begin{onecolentry}
\subsection{Languages}
\begin{highlights}
\item \textbf{Scripting:} TypeScript/JavaScript, Python, Bash
\item \textbf{Programming:} Go, C
\item \textbf{Markup:} HTML/CSS, Markdown, \LaTeX
\end{highlights}

\subsection{Technologies \& Tools}
\begin{highlights}
\item \textbf{Frameworks:} React, Next.js, Node.js, React Native
\item \textbf{Databases:} PostgreSQL, MySQL, SQLite, MongoDB
\item \textbf{DevOps:} Docker, Nginx, HAProxy, ufw, Git
\item \textbf{Other:}  Figma
\end{highlights}

\subsection{Development Environment}
\begin{highlights}
\item \textbf{OS:} GNU/Linux
\item \textbf{Editor:} Neovim
\end{highlights}
\end{onecolentry}

\section{General Skills}
\begin{onecolentry}
\subsection{Languages}
\begin{highlights}
\item \textbf{English :} Proficient
\item \textbf{French :} Native
\end{highlights}
\end{onecolentry}

\end{document}
